% CS-473 SPSOC, Practical work 2: The Mandelbrot set
% Compile with: latexmk -pdf report.tex   (or: pdflatex report.tex, twice)
\documentclass[10pt,a4paper]{article}

\usepackage[T1]{fontenc}
\usepackage[utf8]{inputenc}
\usepackage[margin=2.2cm]{geometry}
\usepackage{booktabs}
\usepackage{amsmath}
\usepackage{amssymb}
\usepackage{xcolor}
\usepackage{listings}
\usepackage{graphicx}
\usepackage{caption}
\usepackage{subcaption}
\usepackage{array}
\usepackage[colorlinks=true,linkcolor=black,citecolor=black,urlcolor=blue]{hyperref}

% ===========================================================================
% Placeholder for numbers that still have to be measured on the Gecko5 board.
% Every \tbd in this document MUST be replaced before handing in.
% ===========================================================================
\newcommand{\tbd}{\textcolor{red}{\textbf{??}}}

\lstset{
  basicstyle=\ttfamily\scriptsize,
  keywordstyle=\color{blue!65!black},
  commentstyle=\color{green!40!black}\itshape,
  numberstyle=\tiny\color{gray},
  language=C,
  numbers=left,
  stepnumber=1,
  frame=single,
  framesep=3pt,
  rulecolor=\color{gray!50},
  breaklines=true,
  columns=fullflexible,
  showstringspaces=false,
  xleftmargin=12pt,
}

\newcommand{\Q}[2]{Q#1.#2}

\title{\vspace{-1.2cm}\textbf{CS-473: System Programming for Systems-on-Chip}\\[2mm]
       \large Practical Work 2: The Mandelbrot Set}
\author{%
  Lisa Suzuki \quad (SCIPER: 408471) \\
  Eduard Vlad \quad (SCIPER: 374438)%
}
\date{\today}

\begin{document}
\maketitle
\vspace{-6mm}

% ===========================================================================
\section{Introduction}
% ===========================================================================

This report documents the implementation of two precision data types:
\begin{enumerate}\itemsep0pt
  \item \textbf{\texttt{fractal\_fxpt}}, a signed \Q{8}{24} fixed-point version;
  \item \textbf{\texttt{fractal\_myflpt}}, a custom 32-bit floating-point format with
        inlined arithmetic.
\end{enumerate}

Both reach the same image as the reference, with $\le 0.03\,\%$ of pixels differing, while
eliminating all calls into libgcc's soft-float library.

\paragraph{Experimental setup.} 
Table~\ref{tab:setup} reports the configuration under which we collected the figures below.
Changes between implementations only involve changes in the number formats.

\begin{table}[h]
\centering\small
\caption{Experimental setup.}
\label{tab:setup}
\begin{tabular}{@{}ll@{}}
\toprule
\multicolumn{2}{@{}l}{\textbf{Platform}}\\
\midrule
CPU                      & OR1300 (OpenRISC~1000 ISA, 5-stage pipeline, no FPU) \\
Instruction cache        & 8\,kiB, direct mapped, FIFO replacement \\
Data cache               & 8\,kiB, 4-way set associative, LRU, write-back \\
Toolchain                & \texttt{or1k-elf-gcc 15.2.0} \\
Compiler flags           & \texttt{-Os -DNDEBUG -D\_\_OR1300\_\_} (project default) \\
\midrule
\multicolumn{2}{@{}l}{\textbf{Workload}}\\
\midrule
Image size               & $512 \times 512$ pixels, \texttt{rgb565} frame buffer \\
Iteration limit          & $n_{\max} = 64$ \\
Viewport                 & $c_x \in [-2, 1)$, \; $c_y \in [-1.5, 1.5)$ \\
Coordinate step          & $\Delta c = 3/512 = 0.005859375$ \\
Total iterations         & $\approx 3.87 \times 10^{6}$ (Section~\ref{sec:verif}) \\
\midrule
\multicolumn{2}{@{}l}{\textbf{Origin of the reported figures}}\\
\midrule
Static code metrics      & \texttt{or1k-elf-size} and \texttt{or1k-elf-objdump} on the ELF \\
Run-time metrics         & OR1300 profiling counters (Section~\ref{sec:board}) \\
Numerical accuracy       & host-side replay against IEEE-754 \texttt{double} (Section~\ref{sec:verif}) \\
\bottomrule
\end{tabular}
\end{table}

% ===========================================================================
\section{Fixed-point implementation}
% ===========================================================================

\subsection{Format definition}
\label{sec:qformat}

Choosing $m$ and $f$ requires a bound on the largest intermediate value, which we derive here for
this the Mandelbrot set algorithm and viewport with $c_x \in [-2,1)$ and $c_y \in [-1.5,1.5]$.
We compute the bounds by bounding every iteration and the trivial step at initialization.

\textit{Iteration Bounds}: Let $x_n, y_n$ be the values at the start of iteration $n$. 
The loop body computes $xx = x_n^2$, $yy = y_n^2$, updates $x$ and $y$ by the Mandelbrot recurrence, and only \emph{then} tests
$xx + yy < 4$. 
Iteration $n+1$ is reached only if $x_n^2 + y_n^2 < 4$. Under that
hypothesis, and using the maximum values $|c_x| \le 2$, $|c_y| \le 1.5$ for our viewport:
\begin{align}
    |x_{n+1}| &= |x_n^2 - y_n^2 + c_x| \;\le\; (x_n^2 + y_n^2) + |c_x| \;<\; 4  + 2 = 6, \\
  |y_{n+1}| &= |2 x_n y_n + c_y|     \;\le\; (x_n^2 + y_n^2) + |c_y| \;<\; 4 + 1.5 = 5.5,
\end{align}
where the second line uses $|2xy| \le x^2 + y^2$. 
The largest value ever written to \texttt{xx + yy} is bounded by
\begin{equation}
  x_{n+1}^2 + y_{n+1}^2 \;<\; 6^2 + 5.5^2 \;=\; 66.25 .
  \label{eq:bound}
\end{equation}

\textit{First Iteration}: The first iteration where $x_1 = c_x$ and $y_1 = c_y$, yields at most $4 + 2.25 = 6.25$ and is
not critical. The format must therefore represent magnitudes up to $66.25$, so it needs
$\lceil \log_2 66.25 \rceil = 7$ integer bits plus one sign bit. Consequently, the integer digit can be represented by $m = 8$ bits leaving 24  fractional bits:
\[
  \boxed{\texttt{fxpt} = \Q{8}{24}, \qquad \text{range } [-128,\,128), \qquad
  \text{resolution } 2^{-24} \approx 5.96 \times 10^{-8}.}
\]

We show empirically, that the value range is not exceeded: the largest value observed for
\texttt{xx + yy} is $34.925$ without overflow events in $3.87 \times 10^6$ iterations.

We wrap literal definitions in a C-macro (\texttt{FX()}) which is constant-folded by the compiler at compile-time. 
The declaration keeps the scale factor configurable and is used for all fractional literals in the program:

\begin{lstlisting}
#define FXPT_FRAC 24                            // In Q(32-FXPT_FRAC).FXPT_FRAC
#define FX(x) ((fxpt)((x) * (1 << FXPT_FRAC)))  // compile-time float -> fxpt conversion
typedef int32_t (fxpt);
\end{lstlisting}

\subsection{Arithmetic operations}

Addition and subtraction use the native 32-bit integer addition, as both operands have the 
same scale. However, multiplication required additional consideration, since the resulting 
\Q{16}{48} product of two \Q{8}{24} fixed-point decimals does not fit in 32 bits:

\begin{lstlisting}
static inline fxpt mul(fxpt a, fxpt b) {
    return (fxpt)(((int64_t) a * b) >> FXPT_FRAC);   // Q8.24 x Q8.24 = Q16.48 -> >>24 -> Q8.24
}
\end{lstlisting}

We leverage the fact that the 8 upper bits are always sign extension.
As the maximum absolute value is bounded by 66.25, the upper 8 bits
are not required to represent the resulting integer digit of the multiplication.

Shifting 24 bits rightward truncates toward $-\infty$ (as opposed to C's integer truncation) and introduces a rounding error of at most $2^{24}$.
In Section~\ref{sec:verif}, we show that this rounding error is harmless.

\subsection{Handling Over- and Underflows}

We handle over- and underflows by \textit{omitting them by construction}.
This makes detecting them at runtime obsolete, as we can prove their absence.

\begin{itemize}\itemsep0pt
  \item \textbf{Overflow} is excluded by the bound~\eqref{eq:bound} together with the choice of
        \Q{8}{24}, as shown above. The same argument covers the intermediates:
        $|xx|, |yy| < 36$ and $|2xy| < 66$, all within $[-128, 128)$.
  \item \textbf{Underflow} represents a loss of precision, and does not wrap: a product smaller
        than $2^{-24}$ truncates to zero by shifting 24 bits rightwards.
\end{itemize}

The above guarantee is viewport-specific. A postulated zoom feature would shrink $\Delta c$ but leave the
bound~\eqref{eq:bound} unchanged (it depends only on the escape threshold and on
$|c| \lesssim 2$), so \Q{8}{24} would remain safe against overflow. 
However, we later show that resolution is insufficient under this data format (cf. Section~\ref{sec:discussion}).

% ===========================================================================
\section{Custom floating-point format: \texttt{myflpt}}
% ===========================================================================

\subsection{Format definition}

We adopt a similar design to the suggested layout: a 32-bit word holding a sign bit, a 23-bit
magnitude and an 8-bit exponent in excess-250.

\begin{center}\small
\setlength{\tabcolsep}{4pt}
\begin{tabular}{|c|c|c|}
\hline
\textbf{bit 31} & \textbf{bits 30\,--\,8} & \textbf{bits 7\,--\,0} \\
sign & magnitude/mantissa $m$ (23 bits, leading one \emph{explicit}) & exponent $e$ (excess-250) \\
\hline
\end{tabular}
\end{center}

A non-zero value is interpreted as
\begin{equation}
  v \;=\; (-1)^{s} \cdot \frac{m}{2^{23}} \cdot 2^{\,e-250},
  \qquad \tfrac{1}{2} \le \frac{m}{2^{23}} < 1 ,
\end{equation}
so the mantissa is normalised to the interval $[0, 1)$. 
The all-zero word \texttt{0x00000000} is the unique encoding of zero. Representable
magnitudes span $2^{-250}$ up to just below $2^{7} = 128$.

This choice is more suited for representing the Mandelbrot set than IEEE 754 for the following reasons:
\begin{itemize}\itemsep1pt
  \item \textbf{Exponent in the least significant bits, mantissa above it.} This ordering makes
      two operations cheap. \textit{First}, an unsigned comparison of two same-sign words is already a
        correct magnitude comparison, so \texttt{order()} and \texttt{is\_smaller()} reduce to
        integer compares. \textit{Second}, doubling a value is a single increment of the exponent field.
  \item \textbf{Excess-250 instead of excess-127.} The algorithm only ever handles magnitudes
        below $66.25$, and squaring drives values toward zero. Biasing the exponent so that the
        representable range is $[2^{-250}, 2^{5})$ in line with the maximum proven bounds for the algorithm.

        Magnitudes saturate at $2^5 = 32$. Over the full image, \texttt{mul()} saturates
        24 times, when it squares a value that has already escaped. This is harmless: the escape
        test compares against $4.0$, and a saturated result of $32 > 4$ still fails the test, so
        the iteration terminates exactly as it should. 
  \item \textbf{Saturation and flush-to-zero instead of $\pm\infty$ and NaN.} There is no useful
        reaction to an infinity in this algorithm, and supporting the special encodings would add a
        test to every operation. Exponent overflow therefore returns the largest representable
        magnitude ($128$, encoded \texttt{0x7FFFFFFF}) and exponent underflow returns zero.
        Section~\ref{sec:myflpt-limits} shows why saturating at 128 cannot corrupt the escape test.
\end{itemize}

\subsection{Arithmetic operations}

All operators are implemented in \texttt{include/myflpt.h} and are declared \texttt{static inline} so that
they can be inlined into the hot loop. 
Section~\ref{sec:inlining} shows that this requires an explicit compiler flag to force inlining.

\paragraph{\texttt{float\_to\_myflpt()}} converts an IEEE-754 \texttt{float} by re-biasing the
exponent, since $1.m \cdot 2^{E-127} = 0.1m \cdot 2^{(E+124)-250}$, and by shifting the 24-bit
significand $1.m$ right by one to place the leading one at bit~22 of the magnitude field. It is
used only on compile-time constants and is fully folded away.

\paragraph{\texttt{add()} and \texttt{sub()}} align, operate and normalise. \texttt{order()} first swaps the operands so that $|a| \ge |b|$, which guarantees that
the exponent difference $d = e_a - e_b$ is non-negative and that the subtraction of magnitudes
cannot go negative. The smaller magnitude is then shifted right by $d$ and the two are added
(equal signs) or subtracted (opposite signs), after which the result is re-normalised. Both
mantissas are widened by \texttt{MYFLPT\_G}~$=8$ guard bits before alignment, so that bits shifted
out during alignment still contribute to the result; the guard bits are dropped at the very end.
The guard $d > 23 + \mathtt{MYFLPT\_G}$ both short-circuits the case where $b$ is negligible and
avoids an undefined shift by more than 31. Subtraction is implemented as
$a - b = a + (-b)$ by flipping the sign bit, which costs one XOR.

\paragraph{\texttt{mul()}} adds the exponents ($e = e_a + e_b - 250$, the extra $-250$ undoing the
double bias), XORs the signs, and forms the $46$-bit product of the two 23-bit magnitudes in a
\texttt{uint64\_t}. Since both magnitudes lie in $[0.5, 1)$, the product lies in $[0.25, 1)$ and
needs at most a single normalisation step, selected by one comparison against $2^{45}$.

\paragraph{\texttt{is\_smaller()}} performs the escape test. Operands with different signs are
ordered by sign alone; same-sign operands are ordered by the unsigned comparison of their
(exponent, magnitude) fields, with the order reversed for two negative values.

\paragraph{Computing $2xy$.} 
This computation deserves a remark: we first attempted to add 1 to the exponent, as it 
occupies the lower bits of the word.
However, it is unsafe: when \texttt{mul()} saturates, the exponent field is already \texttt{0xFF} 
and the increment carries into the magnitude and sign fields, producing \texttt{0x80000000}. 
We therefore use \texttt{two\_xy = add(xy, xy)} instead, which saturates correctly. We verified that both forms
produce a bit-identical image at this zoom level (0 pixels differ), because the
saturating case only ever occurs on the iteration that escapes; the safe form was kept because the
property does not survive a change of viewport.

\subsection{Limitations}
\label{sec:myflpt-limits}

We do not implement subnormals, infinities or NaNs. 
\texttt{float\_to\_myflpt()} maps a subnormal input to zero. 
This is deliberate and cannot be reached from constant inputs.

% ===========================================================================
\section{Verification}
\label{sec:verif}
% ===========================================================================

We replay both fixed and floating point libraries on our host against the IEEE-754 \texttt{double} reference (\texttt{report/verify.c}), using the exact same source headers that are compiled for the OR1300. 
We checked three properties: the accuracy of the individual operators, the ordering
predicate, and the iteration count of every pixel of the final image.

\begin{table}[h]
\centering\small
\caption{Worst-case relative error of the \texttt{myflpt} operators over $10^6$ random operand
pairs drawn from $[-4, 4]$.}
\label{tab:accuracy}
\begin{tabular}{lcc}
\toprule
\textbf{Operation} & \textbf{worst relative error} & \textbf{$\approx$} \\
\midrule
\texttt{float\_to\_myflpt()} & $1.192 \times 10^{-7}$ & $2^{-23}$ \\
\texttt{add()}               & $2.360 \times 10^{-7}$ & $2^{-22}$ \\
\texttt{sub()}               & $2.353 \times 10^{-7}$ & $2^{-22}$ \\
\texttt{mul()}               & $2.380 \times 10^{-7}$ & $2^{-22}$ \\
\midrule
IEEE-754 \texttt{float} (for reference) & $1.192 \times 10^{-7}$ & $2^{-23}$ \\
\bottomrule
\end{tabular}
\end{table}

The bound of $2^{-22}$ for the arithmetic operators is exactly what the format predicts: the result
magnitude can be as small as $0.5$, so one truncated unit of precision $2^{-23}$ represents a relative error
of up to $2^{-22}$. \texttt{is\_smaller()} was checked exhaustively over all $81$ ordered pairs
drawn from $\{-8, -4, -1.5, -0.25, 0, 0.25, 1.5, 4, 8\}$ and agrees with the \texttt{double}
ordering on all of them.

\begin{table}[h]
\centering\small
\caption{Agreement of the full $512\times512$ image with a \texttt{double} reference evaluated at
the same (decoded) coordinates. ``Delta'' is the difference in the returned iteration count.}
\label{tab:image}
\begin{tabular}{lccc}
\toprule
\textbf{Version} & \textbf{pixels differing} & \textbf{share} & \textbf{max.\ delta} \\
\midrule
\texttt{fxpt} (\Q{8}{24}) & $39 / 262\,144$ & $0.0149\,\%$ & 7 \\
\texttt{myflpt}           & $66 / 262\,144$ & $0.0252\,\%$ & 19 \\
\bottomrule
\end{tabular}
\end{table}

Both versions are therefore visually indistinguishable from the reference. The differing pixels all
lie on the fractal boundary, where the iteration is chaotic and an error of one unit of precision in the input
changes the escape iteration. The total work is essentially identical in all three versions
($3\,873\,537$ iterations for \texttt{fxpt}, $3\,873\,492$ for \texttt{myflpt}, $3\,873\,534$ for
the \texttt{double} reference, a spread of $0.002\,\%$).

% ===========================================================================
\section{Performance}
\label{sec:perf}
% ===========================================================================

The four programs (float, fxpt, myflpt, myflpt\_always\_inline) were instrumented with the OR1300 profiling module
(\texttt{support/include/perf.h}). 

\paragraph{Methodology.} Each configuration was measured three times, with a board reset between
runs to start with cold instruction and data caches.
We report cycle counts rather than the \texttt{hh:mm:ss.mmm}.

The executed-instruction counter is deterministic for a given binary. 
The run-time, stall and bus-idle counters do vary slightly, because two bus masters are asynchronous to the program: the SDRAM controller issues
an auto-refresh every $\sim 78\,\mu$s irrespective of CPU activity.

Both loads are \emph{common-mode}, identical in all four binaries, so they shift every row of Table~\ref{tab:board} equally and do not distort the comparison. 

\begin{table}[h]
\centering\small
\caption{Measured on the Gecko5. Median of three runs; raw UART logs in
\texttt{report/measurement/}. The derived clock is $5\,971\,036\,224/100.524 = 59.40$\,MHz and is
identical for all four binaries.}
\label{tab:board}
\begin{tabular}{lrrrrrr}
\toprule
\textbf{Version} & \textbf{\begin{tabular}{@{}c@{}}run time\\{}[s]\end{tabular}}
& \textbf{\begin{tabular}{@{}c@{}}run time\\{}[cycles]\end{tabular}}
& \textbf{\begin{tabular}{@{}c@{}}executed\\instructions\end{tabular}}
& \textbf{\begin{tabular}{@{}c@{}}instr.\ per\\iteration\end{tabular}}
& \textbf{CPI} & \textbf{speed-up} \\
\midrule
\texttt{fractal\_flpt}            & $100.524$ & $5\,971\,036\,224$ & $3\,292\,407\,183$ & $850.0$ & $1.814$ & $1.00\times$ \\
\texttt{fractal\_fxpt}            & $8.853$   & $525\,909\,811$     & $479\,408\,457$     & $123.8$ & $1.097$ & $\mathbf{11.35\times}$ \\
\texttt{fractal\_myflpt}          & $39.391$  & $2\,339\,826\,624$ & $2\,039\,215\,915$ & $526.5$ & $1.147$ & $\mathbf{2.55\times}$ \\
\texttt{..\_always\_inline}        & $30.738$  & $1\,825\,845\,296$ & $1\,632\,802\,044$ & $421.5$ & $1.118$ & $\mathbf{3.27\times}$ \\
\bottomrule
\end{tabular}
\end{table}

\begin{table}[h]
\centering\small
\caption{Memory-system counters, median of three runs. Bus occupancy is
$1-(\text{bus idle}/\text{run time})$.}
\label{tab:memsys}
\begin{tabular}{lrrrrr}
\toprule
\textbf{Version} & \textbf{I\$ misses} & \textbf{\begin{tabular}{@{}c@{}}I\$ miss penalty\\{}[cycles]\end{tabular}}
& \textbf{\begin{tabular}{@{}c@{}}\% of\\run time\end{tabular}}
& \textbf{D\$ misses} & \textbf{\begin{tabular}{@{}c@{}}bus\\occupancy\end{tabular}} \\
\midrule
\texttt{fractal\_flpt}            & $24\,277\,779$ & $1\,647\,278\,253$ & $\mathbf{27.6\,\%}$   & $16\,400$ & $40.1\,\%$ \\
\texttt{fractal\_fxpt}            & $32$            & $1\,586$             & $<0.001\,\%$          & $16\,389$ & $28.2\,\%$ \\
\texttt{fractal\_myflpt}          & $52$            & $3\,506$             & $<0.001\,\%$          & $16\,389$ & $27.9\,\%$ \\
\texttt{..\_always\_inline}        & $112$           & $7\,685$             & $<0.001\,\%$          & $16\,390$ & $27.9\,\%$ \\
\bottomrule
\end{tabular}
\end{table}

\subsection{Always Inline}
\label{sec:inlining}

Declaring custom operators \texttt{static inline} turns out \emph{not} to be sufficient at the project's default optimisation level: at \texttt{-Os} as gcc keeps out-of-line copies of \texttt{add()} and \texttt{mul()} and calls them. 

We therefore built a \texttt{fractal\_myflpt\_always\_inline}, which
differs from \texttt{fractal\_myflpt} by placing \texttt{\_\_attribute\_\_((always\_inline))} to function definitions.

\begin{table}[h]
\centering\small
\caption{Effect of forcing inlining of the \texttt{myflpt} operators (measured, median of three runs).}
\label{tab:inline}
\begin{tabular}{lrrrrr}
\toprule
\textbf{Build} & \textbf{.text [B]} & \textbf{calls / iter.} &
\textbf{\begin{tabular}{@{}c@{}}instr.\ per\\iteration\end{tabular}} &
\textbf{run time [s]} & \textbf{I\$ misses} \\
\midrule
\texttt{static inline} only (\texttt{-Os})    & $18\,236$ & 8 & $526.5$ & $39.391$ & 52 \\
\texttt{+ always\_inline}                    & $20\,216$ & 3 & $421.5$ & $\mathbf{30.738}$ & 112 \\
\midrule
\emph{change}                                & $+10.9\,\%$ & $-5$ & $-104.9$ & $-22.0\,\%$ & $+60$ \\
\bottomrule
\end{tabular}
\end{table}

Forcing inlining removes all eight library calls, leaving only the three \texttt{\_\_muldi3} calls
that the 64-bit mantissa product requires, the same number the fixed-point version pays. The
program runs $1.28\times$ faster.

% ===========================================================================
\section{Discussion}
\label{sec:discussion}
% ===========================================================================

\paragraph{Using Fixed or Floating Point.} For this viewport the fixed-point version wins clearly: $11.35\times$ against $3.27\times$ for the best custom-float build, a smaller binary, and
marginally better accuracy (Table~\ref{tab:image}), because its $2^{-24}$ absolute resolution is
finer than \texttt{myflpt}'s $2^{-22}$ relative resolution for values of order one, which is all
this viewport contains. 

However, with a zoom feature $\Delta c$ shrinks while the escape bound stays at
$4.0$; \Q{8}{24} keeps its 24 fractional bits regardless and becomes visibly quantised after $2^{20}$ magnification, whereas \texttt{myflpt} redistributes precision automatically. 

% ===========================================================================
\section{Conclusion}
% ===========================================================================

On the board, the fixed-point version runs $\mathbf{11.35\times}$ faster than the library
\texttt{float} and the custom floating-point version $\mathbf{2.55\times}$ faster, rising to
$\mathbf{3.27\times}$ once the operators are always inlined.

\end{document}
